\section{Discussion and Future Work}
While the preliminary results demonstrate the feasibility of the proposed architecture, a comprehensive evaluation requires systematic comparison with established forecasting models. Future work will benchmark the NS-SDN architecture against classical econometric methods such as ARIMA and vector autoregressions, as well as neural forecasting models including recurrent networks and feedforward architectures.
Forecast comparisons will be conducted using rolling-origin evaluation procedures commonly employed in forecasting studies \cite{west1996predictive}. Statistical significance of forecast improvements will be assessed using tests for equal predictive accuracy such as the Diebold-Mariano test \cite{diebold1995accuracy}.
Beyond predictive performance, the interpretability of the learned spectral components provides an additional avenue for analysis. Examining the evolution of amplitude and frequency parameters may offer insight into changes in economic cycle dynamics or regime transitions.
These extensions will form the basis of a larger empirical study evaluating the effectiveness of state-driven spectral decomposition models for econometric forecasting.
